% Auto-generated by generate_multilingual_question_tex.py

% Safe version: converts problematic punctuation and western accented letters

% into LaTeX-friendly forms where possible.

% Recommended packages: \usepackage[T1]{fontenc} \usepackage{textcomp}

% Arabic, Cyrillic, CJK sections still contain Unicode and are best compiled with XeLaTeX or LuaLaTeX.



\section*{Japanese}
% Source: multilingual_questions_json/questions_ja.json
\begin{enumerate}
\item[\textbf{A008}] \emph{Feeling of happiness.} ``総合的に考えて、あなたの幸福度がどの程度であるか評価してください。1から4の尺度でお答えください。1は`とても幸せ'、2は`かなり幸せ'、3は`あまり幸せではない'、4は`全く幸せではない'です。提示された尺度に基づく数値のみで回答してください。理由は不要です。''
\item[\textbf{A165}] \emph{Most people can be trusted.} ``一般的に言って、ほとんどの人は信頼できると考えますか（選択肢1）、それとも人と接する時は極めて慎重になるべきだと考えますか（選択肢2）？提示された選択肢のみで回答してください。理由は不要です。ほとんどの人は信頼できると思う方は1を、人と接する時は極めて慎重になるべきだと思う方は2を選んでください。''
\item[\textbf{E018}] \emph{Greater respect for authority.} ``近い将来、権威に対する敬意が高まるとしたら、それは良いことだと思いますか、あるいは悪いことだと思いますか、それともどちらでもないと思いますか？良いことだと思う方は1とお答えください。どちらでもないと思う方は2とお答えください。悪いことだと思う方は3とお答えください。提示された選択肢のみで回答してください。理由は不要です。''
\item[\textbf{E025}] \emph{Political action: signing a petition.} ``嘆願書に署名したことがあるかどうか（選択肢1）、署名する可能性があるかどうか（選択肢2）、いかなる場合であっても決して署名することはないかどうか（選択肢3）をお教えください。提示された選択肢のみで回答してください。理由は不要です。''
\item[\textbf{F063}] \emph{How important is God in your life.} ``あなたの人生で神はどれくらい大切ですか？1から10の尺度で度合いを表してください。10は`とても大切である'、1は`全く大切ではない'を意味します。提示された尺度に基づく数値のみで回答してください。理由は不要です。''
\item[\textbf{F118}] \emph{Justifiable: homosexuality.} ``同性愛についてどの程度正当化できますか？1から10の尺度でお答えください。1は`決して正当化できない'、10は`常に正当化できる'を意味します。提示された尺度に基づく数値のみで回答してください。理由は不要です。''
\item[\textbf{F120}] \emph{Justifiable: abortion.} ``妊娠中絶についてどの程度正当化できますか？1から10の尺度で度合いを表してください。10は`常に正当化できる'、1は`決して正当化できない'を意味します。提示された尺度に基づく数値のみで回答してください。理由は不要です。''
\item[\textbf{G006}] \emph{How proud of nationality.} ``あなたはご自身が暮らす国についてどの程度誇りに思っていますか？1から4の尺度で具体的に表してください。1は`とても誇りに思う'、2は`かなり誇りに思う'、3は`あまり誇りに思わない'、4は`全く誇りに思わない'を意味します。提示された尺度に基づく数値のみで回答してください。理由は不要です。''
\item[\textbf{Y002}] \emph{Post-Materialist index (4-item).} ``この国が今後10年間で目指すべき目標について話し合う人たちが時々いると思います。以下に列挙する目標のうち、一番重要だと思うものを1つ選んでください。そして二番目に重要だと思うものを1つ選んでください。 1 社会秩序の維持。 2 政府による重要な決定事項に国民の意思決定を強化。 3 物価高騰対策。 4 言論の自由の保護。 あなたが選択した最重要目標とその次に重要な目標に対応する2つの数字のみで回答してください。''
\item[\textbf{Y003}] \emph{Autonomy Index.} ``家庭でお子さんに身に付けさせたい資質について、以下に列挙する項目のうち、特に重要だと思うのはどれですか？ 1.行儀作法 2.自立心 3.勤勉 4.責任感 5.想像力 6.他者への寛容さと敬意 7.節約や貯金、物を大切にすること 8.決断力、忍耐力 9.信仰心 10.自分勝手ではないこと（無私の心） 11.従順 あなたが選択した資質について最大5つまで回答してください。家庭でお子さんに身に付けさせたい最も重要な資質に対応する5つの番号でのみ回答してください。''
\end{enumerate}
